\scanpage{223}
Cette fonction est déterminée, indépendamment du choix de $C_s$, modulo une expression de la forme
\[
(3)\qquad \sum D_{is}\int_{Z_s}\omega_{is}\quad
\deleted{\unsure{=\int_{Z_s}\sum D_{is}\omega_{is}}}
\]
\deleted{(\uline{NB} Le fait que $Z_s$ est un \uline{1-cycle} permet de}

Supposons que les $D_i,\omega_i$ soient tels que l'expression précédente soit nulle, quelle que soit \deleted{\unsure{et}} la famille continue $s\longrightarrow Z_s$ de 1-cycles \deleted{$\unsure{Z}$} sur un ouvert de $S$. \deleted{Alors l'expression} \inserted{On dira} que \deleted{\unsure{les}} $\mu\unsure{=|:}(D_i,\omega_i)_{i\in I}$ \uline{définit une équation de Picard-Fuchs sur} $X/S$. Alors pour deux sections hol. \unsure{quelc.} $a,b$ de $X/S$, l'expression
\[
(4)\qquad \mu(a,b)=\sum_i D_{is}\int_{a(s)}^{b(s)}\omega_{is}
\]
est bien déterminée, et est une fonction holomorphe sur $S$. Bien sûr on aura
\[
(5)\qquad \mu(a,b)+\mu(b,c)=\mu(a,c).
\]

Cherchons une condition suffisante maniable, \inserted{\unsure{d'une forme diff.|d'une forme différentielle}} \deleted{pour} \deleted{un \unsure{critère}\inserted{\unsure{d'allure|faible}} \uuline{infinitésimal} et en plus global,} pour que $\mu$ soit de Pic. Fuchs. Notons
\margin{222}
